\scanpage{38}
% archive foliation: 37
Il y a lieu de \unsure{linéariser|généraliser}
$1^\circ)$ au cas où on a un anneau
de base $k$ de Dedekind,
et $2^\circ)$ un $k$ quelconque \inserted{(\unsure{evt. commut.|est commut.})}
(donc $C\subset\widehat\Delta_{k\text{-}mod}$). La \unsure{double|seule}
conclusion à retenir, c'est
qu'on connaît les $\operatorname{Hom}$ de
$X_*$ dans les $k$-modules str.
$\partial$-\deleted{\unsure{parf.}}\inserted{\unsure{parfaits}} \inserted{$M$ (comme foncteur en $M$)} dès qu'on connaît
le complexe \inserted{de cochaînes ``non dég''} $C^{*!}(X_*,k)$ à coeff. dans
$k$. Notre propos principal
est de voir comment reconstituer
les structures multiplicatives
et ``puissances divisées'' sur ce
foncteur
\[
M\longmapsto F_{X_*}(M_*)=F(M_*)=\operatorname{Hom}_{\unsure{\widehat\Delta|S}}(X_*,M_*)
\]
i.e. les lois
\[
\left\{\begin{aligned}
F(M_*)\otimes_k F(N_*)&\longrightarrow F(M_*\otimes N_*)\\
\Gamma^iF(\deleted{\unsure{\Gamma^i}}M_*)&\longrightarrow F(\Gamma^i(M_*))
\end{aligned}\right.
\]
en termes de \deleted{\unsure{la}} structures
sur le complexe \inserted{de cochaînes} $C^*=C^{*!}(X_*,k)$.
Il faudrait examiner si pour
$k$ de car $0$ \inserted{$C^*$ \unsure{avec}} cette structure suppl.
